\documentclass[11pt]{article}

% ===================== PACKAGES =====================
\usepackage[a4paper,margin=1in]{geometry}
\usepackage{amsmath,amssymb}
\usepackage{enumitem}
\usepackage{fancyhdr}
\usepackage{xcolor}
\usepackage{setspace}

% ===================== HEADER & FOOTER =====================
\pagestyle{fancy}
\fancyhf{}
\lhead{CSE 400: Fundamentals of Probability in Computing}
\rhead{Lecture 6 Scribe}
\cfoot{\thepage}

% ===================== TITLE =====================
\title{
    \normalsize School of Engineering and Applied Sciences, Ahmedabad University \\
    \vspace{0.2cm}
    \textbf{CSE 400: Fundamentals of Probability in Computing} \\
    \vspace{0.2cm}
    \Large Lecture 6: Discrete RVs, Expectation and Problem Solving
}

\date{}

\begin{document}
\maketitle
\onehalfspacing

% ==========================================================
\section{Previous Lecture Recap: Random Variables (RVs)}

\subsection{Motivation and Concept of a Random Variable}

A \textbf{random variable} is a variable whose value depends on the outcome of a random experiment.

From the lecture:
\begin{itemize}
    \item The \textbf{distribution of a random variable} can be visualized as a \textbf{bar diagram}.
    \item The \textbf{x-axis} represents the values that the random variable can take.
    \item The \textbf{height of the bar} at value $a_i$ represents
    \[
        \Pr[X = a_i]
    \]
    \item Each probability is computed by finding the probability of the \textbf{corresponding event in the sample space}.
\end{itemize}

% ==========================================================
\subsection{Random Variable Example}

\textbf{Example 1: Tossing 3 Fair Coins}

Let:
\begin{itemize}
    \item The experiment consist of tossing \textbf{3 fair coins}
    \item $Y$ = number of heads observed
\end{itemize}

Possible values of $Y$:
\[
Y \in \{0,1,2,3\}
\]

Step-by-step computation of probabilities:
\begin{itemize}
    \item $\Pr(Y = 0) = \Pr(t,t,t) = \frac{1}{8}$
    \item $\Pr(Y = 1) = \Pr(t,t,h), \Pr(t,h,t), \Pr(h,t,t) = \frac{3}{8}$
    \item $\Pr(Y = 2) = \Pr(h,h,t), \Pr(h,t,h), \Pr(t,h,h) = \frac{3}{8}$
    \item $\Pr(Y = 3) = \Pr(h,h,h) = \frac{1}{8}$
\end{itemize}

Verification:
\[
\sum_{i=0}^{3} \Pr(Y=i) = \frac{1}{8}+\frac{3}{8}+\frac{3}{8}+\frac{1}{8} = 1
\]

% ==========================================================
\section{Probability Mass Function (PMF)}

\subsection{Concept}

A random variable that can take \textbf{at most a countable number of values} is called a \textbf{discrete random variable}.

Let:
\begin{itemize}
    \item $X$ be a discrete random variable
    \item Range:
    \[
    R_X = \{x_1, x_2, x_3, \dots\}
    \]
\end{itemize}

The function:
\[
P_X(x_k) = \Pr(X = x_k), \quad k = 1,2,3,\dots
\]
is called the \textbf{Probability Mass Function (PMF)} of $X$.

Property:
\[
\sum_{k=1}^{\infty} P_X(x_k) = 1
\]

% ==========================================================
\subsection{PMF Example (Given in Lecture)}

The PMF of a random variable $X$ is given by:
\[
p(i) = c \frac{\lambda^i}{i!}, \quad i = 0,1,2,\dots
\]
where $\lambda > 0$.

\subsubsection*{Step 1: Use normalization condition}

\[
\sum_{i=0}^{\infty} p(i) = 1
\]

\[
c \sum_{i=0}^{\infty} \frac{\lambda^i}{i!} = 1
\]

Using:
\[
e^\lambda = \sum_{i=0}^{\infty} \frac{\lambda^i}{i!}
\]

So:
\[
c e^\lambda = 1
\Rightarrow c = e^{-\lambda}
\]

\subsubsection*{Step 2: Compute $\Pr[X = 0]$}

\[
\Pr[X=0] = e^{-\lambda} \frac{\lambda^0}{0!} = e^{-\lambda}
\]

\subsubsection*{Step 3: Compute $\Pr[X > 2]$}

\[
\Pr[X>2] = 1 - \Pr[X \le 2]
\]

\[
= 1 - \left( \Pr[X=0] + \Pr[X=1] + \Pr[X=2] \right)
\]

\[
= 1 - \left( e^{-\lambda} + \lambda e^{-\lambda} + \frac{\lambda^2}{2} e^{-\lambda} \right)
\]

% ==========================================================
\section{Independent Events}

\subsection{Definition}

Two events $A$ and $B$ are \textbf{independent} if:
\[
\Pr(A \cap B) = \Pr(A)\Pr(B)
\]

% ==========================================================
\section{Bayes’ Theorem (Recap)}

\subsection{Bayes’ Formula}

Using:
\[
\Pr(A \cap B) = \Pr(B|A)\Pr(A)
\]

We obtain:
\[
\Pr(B_i|A) = \frac{\Pr(A|B_i)\Pr(B_i)}{\sum_{j=1}^{n} \Pr(A|B_j)\Pr(B_j)}
\]

This is known as \textbf{Bayes’ Formula (Proposition 3.1)}.

Definitions:
\begin{itemize}
    \item $\Pr(B_i)$: a priori probability
    \item $\Pr(B_i|A)$: posterior probability
\end{itemize}

% ==========================================================
\subsection{Bayes’ Theorem Example 1: Auditorium with 30 Rows}

Given:
\begin{itemize}
    \item 30 rows
    \item Row 1 has 11 seats
    \item Row 2 has 12 seats
    \item \dots
    \item Row 30 has 40 seats
\end{itemize}

Procedure:
\begin{enumerate}
    \item A row is selected uniformly at random
    \item A seat is selected uniformly within the chosen row
\end{enumerate}

Compute:
\begin{itemize}
    \item $\Pr(\text{Seat 15} \mid \text{Row 20})$
    \item $\Pr(\text{Row 20} \mid \text{Seat 15})$
\end{itemize}

\[
\Pr(S_{15} | R_{20}) = \frac{1}{30}
\]

\[
\Pr(S_{15}) = \sum_{k=5}^{30} \Pr(S_{15} | R_k)\Pr(R_k)
\]

\[
= \sum_{k=5}^{30} \frac{1}{k+10} \cdot \frac{1}{30}
\]

\[
\Pr(R_{20}|S_{15}) = \frac{\Pr(S_{15}|R_{20})\Pr(R_{20})}{\Pr(S_{15})}
\]

% ==========================================================
\section{Types of Discrete Random Variables}

\begin{itemize}
    \item \textbf{Bernoulli Random Variable}: Single trial, two outcomes
    \item \textbf{Binomial Random Variable}: Fixed number of trials, independent trials, same probability of success
    \item \textbf{Geometric Random Variable}: Counts trials until first success
    \item \textbf{Poisson Random Variable}: Models count of events in fixed interval
\end{itemize}

% ==========================================================
\section{Expectation of Random Variables}

\subsection{Definition}

\[
\mathbb{E}[X] = \sum_x x \Pr(X=x)
\]

\subsection{Expectation of a Function of a Random Variable}

\[
\mathbb{E}[g(X)] = \sum_x g(x)\Pr(X=x)
\]

\subsection{Linearity of Expectation}

\[
\mathbb{E}[aX+b] = a\mathbb{E}[X] + b
\]

\subsection{Moments and Central Moments}

\begin{itemize}
    \item Variance:
    \[
    \text{Var}(X) = \mathbb{E}[(X-\mu)^2]
    \]
    \item Skewness: third central moment
    \item Kurtosis: fourth central moment
\end{itemize}

% ==========================================================
\section{Cumulative Distribution Function (CDF)}

\subsection{Definition}

\[
F_X(x) = \Pr(X \le x)
\]

\subsection{Properties}

\begin{itemize}
    \item Non-decreasing
    \item $\lim_{x\to -\infty} F_X(x)=0$
    \item $\lim_{x\to \infty} F_X(x)=1$
\end{itemize}

% ==========================================================
\section{Probability Density Function (PDF)}

\subsection{Definition}

\[
f_X(x) = \frac{d}{dx} F_X(x)
\]

\subsection{Property}

\[
\int_{-\infty}^{\infty} f_X(x)\,dx = 1
\]

\vfill
\begin{center}
\textbf{End of Lecture 6 Scribe — Strictly Aligned with Course Slides}
\end{center}

\end{document}